\documentclass[a4paper,10pt]{article}
\usepackage[utf8]{inputenc}
\usepackage[T1]{fontenc}
\usepackage[french]{babel}
\usepackage{amsmath,amssymb}
\usepackage{geometry}
\usepackage{array,tabularx}
\usepackage{tikz}
\geometry{margin=1.6cm}
\pagestyle{empty}
\newcommand{\ligne}{\par\vspace{0.55cm}\noindent\dotfill\par}
\newcommand{\carreaux}[1]{\begin{center}\begin{tikzpicture}\draw[step=0.5cm,gray!35,very thin] (0,0) grid (17,#1);\end{tikzpicture}\end{center}}
\newenvironment{exercice}[2]{\paragraph{Exercice #1 \hfill #2 points}\mbox{}\par}{\medskip}
\begin{document}
\begin{center}\Large\bfseries Devoir surveillé 19 - Les translations\end{center}
\vspace{2mm}
\begin{exercice}{1}{4}Décrire une translation à l'aide d'un vecteur ou d'un couple de points.\end{exercice}\begin{exercice}{2}{5}Construire l'image d'un triangle par la translation qui transforme $A$ en $B$.\carreaux{5}\end{exercice}\begin{exercice}{3}{5}Une translation conserve-t-elle les longueurs, les angles et l'alignement ? Répondre et justifier.\end{exercice}\begin{exercice}{4}{6}Sur un quadrillage, le point $M(2;3)$ est transformé en $M'(7;1)$. Donner les coordonnées des images de $N(-1;4)$ et $P(0;0)$ par la même translation.\end{exercice}\end{document}
